% Part: first-order-logic
% Chapter: syntax-and-semantics
% Section: formation-sequences

\documentclass[../../../include/open-logic-section]{subfiles}

\begin{document}

\olfileid[hi]{pl}{syn}{fseq}

\olsection{निर्माण अनुक्रम}

!!{formula}s को आगमनात्मक परिभाषा से परिभाषित करना, और उसकी पूरक विधि के
रूप में आगमन द्वारा !!{formula}s के गुणधर्म सिद्ध करना, सुरुचिपूर्ण और कुशल
दृष्टिकोण है। किंतु !!{formula}s के निर्माण को मूल अवयवों से चरण-दर-चरण
देखना भी उपयोगी हो सकता है; यहाँ हम इसके लिए \emph{निर्माण अनुक्रम} की
धारणा का उपयोग करेंगे।

\begin{defn}[सूत्रों के निर्माण अनुक्रम]
\ollabel{defn:fseq-frm}
परिमित अनुक्रम $\tuple{!A_0,\dotsc,!A_n}$ में प्रत्येक पद भाषा~$\Lang L_0$
के प्रतीकों की शृंखला हो। यह $!A$ का \emph{निर्माण अनुक्रम} है यदि
$!A \ident !A_n$ और प्रत्येक $i \leq n$ के लिए या तो $!A_i$ परमाणु सूत्र
हो, अथवा ऐसे $j,k < i$ हों कि निम्नलिखित में से कोई एक शर्त पूरी हो:
\begin{enumerate}
    \tagitem{prvNot}{$!A_i \ident \lnot !A_j$।}{}%
    \tagitem{prvAnd}{$!A_i \ident (!A_j \land !A_k)$।}{}%
    \tagitem{prvOr}{$!A_i \ident (!A_j \lor !A_k)$।}{}%
    \tagitem{prvIf}{$!A_i \ident (!A_j \lif !A_k)$।}{}%
    \tagitem{prvIff}{$!A_i \ident (!A_j \liff !A_k)$।}{}%
\end{enumerate}
\end{defn}

\begin{ex}
\[
\tuple{
    \Obj p_0,
    \Obj p_1,
    (\Obj p_1 \land \Obj p_0),
    \lnot (\Obj p_1 \land \Obj p_0)
}
\]
यह $\lnot (\Obj p_1 \land \Obj p_0)$ का एक निर्माण अनुक्रम है, और यह भी:
\[
\tuple{
    \Obj p_0,
    \Obj p_1,
    \Obj p_0,
    (\Obj p_1 \land \Obj p_0),
    (\Obj p_0 \lif \Obj p_1),
    \lnot (\Obj p_1 \land \Obj p_0)
}.
\]
%
दूसरे उदाहरण से स्पष्ट है कि निर्माण अनुक्रमों में ``अनावश्यक अंश'' हो सकते
हैं: ऐसे सूत्र जो अतिरिक्त हों या निर्माण में कोई योगदान न दें।
\end{ex}

\begin{prop}
\ollabel{prop:formed}
जो भी !!{formula}~$!A$, $\Frm[L_0]$ में हो, उसका एक निर्माण अनुक्रम है।
\end{prop}

\begin{proof}
मानें कि $!A$ परमाणु है। तब अनुक्रम~$\tuple{!A}$, $!A$ का निर्माण अनुक्रम है।
%
अब मानें कि $!B$ और~$!C$ के निर्माण अनुक्रम क्रमशः
$\tuple{!B_0,\dotsc,!B_n}$ और $\tuple{!C_0,\dotsc,!C_m}$ हैं।
%
\begin{enumerate}
  \tagitem{prvNot}{यदि $!A \ident \lnot !B$ हो, तो
      $\tuple{!B_0,\dotsc,!B_n,\lnot !B_n}$, $!A$ का निर्माण अनुक्रम है।}{}
  \tagitem{prvAnd}{यदि $!A \ident (!B \land !C)$ हो, तो
      $\tuple{!B_0,\dotsc,!B_n,!C_0,\dotsc,!C_m,(!B_n \land !C_m)}$,
      $!A$ का निर्माण अनुक्रम है।}{}
  \tagitem{prvOr}{यदि $!A \ident (!B \lor !C)$ हो, तो
      $\tuple{!B_0,\dotsc,!B_n,!C_0,\dotsc,!C_m,(!B_n \lor !C_m)}$,
      $!A$ का निर्माण अनुक्रम है।}{}
  \tagitem{prvIf}{यदि $!A \ident (!B \lif !C)$ हो, तो
      $\tuple{!B_0,\dotsc,!B_n,!C_0,\dotsc,!C_m,(!B_n \lif !C_m)}$,
      $!A$ का निर्माण अनुक्रम है।}{}
  \tagitem{prvIff}{यदि $!A \ident (!B \liff !C)$ हो, तो
      $\tuple{!B_0,\dotsc,!B_n,!C_0,\dotsc,!C_m,(!B_n \liff !C_m)}$,
      $!A$ का निर्माण अनुक्रम है।}{}
\end{enumerate}
!!{formula}s पर आगमन के सिद्धांत से प्रत्येक !!{formula} का एक निर्माण
अनुक्रम है।
\end{proof}

हम इसका विलोम भी सिद्ध कर सकते हैं। यह महत्त्वपूर्ण है, क्योंकि इससे पता चलता
है कि सूत्रों को परिभाषित करने की हमारी दोनों विधियाँ समतुल्य हैं: वे समान
परिणाम देती हैं। इसका यह भी अर्थ है कि हम निर्माण अनुक्रमों की लंबाई पर साधारण
आगमन का उपयोग करके सूत्रों के विषय में प्रमेय सिद्ध कर सकते हैं।

\begin{lem}
\ollabel{lem:fseq-init}
मानें कि $\tuple{!A_0,\dotsc,!A_n}$, $!A_n$ का निर्माण अनुक्रम है और
$k \leq n$। तब $\tuple{!A_0,\dotsc,!A_k}$, $!A_k$ का निर्माण अनुक्रम है।
\end{lem}

\begin{proof}
अभ्यास।
\end{proof}

\begin{thm}
\ollabel{thm:fseq-frm-equiv}
$\Frm[L_0]$ भाषा~$\Lang L_0$ की उन सभी प्रतीक-शृंखलाओं का समुच्चय है
जिनका कोई निर्माण अनुक्रम है।
\end{thm}

\begin{proof}
$F$ भाषा~$\Lang L_0$ की उन सभी प्रतीक-शृंखलाओं का समुच्चय हो जिनका कोई
निर्माण अनुक्रम है। \olref[pl][syn][fseq]{prop:formed} में हम देख चुके हैं
कि $\Frm[L_0] \subseteq F$, इसलिए अब हम विलोम सिद्ध करते हैं।

मानें कि $!A$ का निर्माण अनुक्रम $\tuple{!A_0,\dotsc,!A_n}$ है। हम सिद्ध
करते हैं कि $!A \in \Frm[L_0]$; इसके लिए~$n$ पर प्रबल आगमन करते हैं। हमारी आगमनात्मक
परिकल्पना है कि लंबाई $m < n$ वाला निर्माण अनुक्रम रखने वाली प्रत्येक
प्रतीक-शृंखला~$\Frm[L_0]$ में है। निर्माण अनुक्रम की परिभाषा से या तो
$!A_n$ परमाणु है, अथवा ऐसे $j,k < n$ अवश्य हैं कि निम्नलिखित में से कोई
एक स्थिति हो:
\begin{enumerate}
    \tagitem{prvNot}{$!A_n \ident \lnot !A_j$।}{}%
    \tagitem{prvAnd}{$!A_n \ident (!A_j \land !A_k)$।}{}%
    \tagitem{prvOr}{$!A_n \ident (!A_j \lor !A_k)$।}{}%
    \tagitem{prvIf}{$!A_n \ident (!A_j \lif !A_k)$।}{}%
    \tagitem{prvIff}{$!A_n \ident (!A_j \liff !A_k)$।}{}%
\end{enumerate}
अब हम स्थितियों के अनुसार तर्क करते हैं। यदि $!A_n$ परमाणु है, तो
$!A_n \in \Frm[L_0]$। अन्यथा मानें कि $!A \equiv
(!A_j \land !A_k)$। \olref[pl][syn][fseq]{lem:fseq-init} से
$\tuple{!A_0,\dotsc,!A_j}$ और $\tuple{!A_0,\dotsc,!A_k}$ क्रमशः $!A_j$
और~$!A_k$ के निर्माण अनुक्रम हैं। चूँकि ये $!A$ के निर्माण अनुक्रम के
उचित आरंभिक उप-अनुक्रम हैं, इन दोनों की लंबाई~$n$ से कम है। अतः आगमनात्मक
परिकल्पना से $!A_j$ और~$!A_k$, $\Frm[L_0]$ में हैं; इसलिए !!a{formula}
की परिभाषा से $(!A_j \land !A_k)$ भी उसमें है। अन्य स्थितियाँ समानांतर
तर्क से सिद्ध होती हैं।
\end{proof}

\end{document}
